% 华为杯统一附录模板。
% 本文件由 build_latex.py 复制到论文目录并在导言区实际 \input；
% 生成内容只能调用这里的接口，不得在章节片段中另设代码样式。
\usepackage{minted}
\usepackage{xcolor}
\usepackage{tcolorbox}
\tcbuselibrary{breakable,skins}
\usepackage{longtable}
\usepackage{booktabs}
\usepackage{array}
\usepackage{needspace}
\usemintedstyle{lovelace}

% 克制的浅色学术代码配色。语法 token 由 Pygments 的 lovelace 浅色主题生成；
% 下列颜色只负责纸面、边线和行号，不覆盖 Pygments token 语义。
\definecolor{AppendixStripeWhite}{HTML}{FFFFFF}
\definecolor{AppendixStripeGray}{HTML}{F8F8F8}
\definecolor{AppendixDivider}{HTML}{63B96C}
\definecolor{AppendixHorizontalRule}{HTML}{555B57}
\definecolor{AppendixLineNumber}{HTML}{6F8175}
\definecolor{AppendixCaption}{HTML}{343A36}
\newfontfamily\AppendixCodeLatin{Consolas}[Scale=MatchLowercase]
\newcommand{\AppendixCodeFont}{\songti\AppendixCodeLatin}
\renewcommand{\myFont}{\AppendixCodeFont}

% Pygments 仍负责词法分类；在 lovelace 浅色基底上将论文要求的
% “关键词紫、字符串绿、数字/类型棕橙、注释灰绿”固化为低饱和调色板。
% Minted 在首次代码块前载入 Pygments style，因此在 MintedVerbatim
% 开始时覆写 token 显示色，避免被后载入的 styledef 反向覆盖。
\makeatletter
\newcommand{\AppendixPygmentsPalette}{%
  \@namedef{PYG@tok@k}{\def\PYG@tc####1{\textcolor[HTML]{8B4E8F}{####1}}}%
  \@namedef{PYG@tok@kc}{\def\PYG@tc####1{\textcolor[HTML]{8B4E8F}{####1}}}%
  \@namedef{PYG@tok@kd}{\def\PYG@tc####1{\textcolor[HTML]{8B4E8F}{####1}}}%
  \@namedef{PYG@tok@kn}{\def\PYG@tc####1{\textcolor[HTML]{8B4E8F}{####1}}}%
  \@namedef{PYG@tok@kp}{\def\PYG@tc####1{\textcolor[HTML]{8B4E8F}{####1}}}%
  \@namedef{PYG@tok@kr}{\def\PYG@tc####1{\textcolor[HTML]{8B4E8F}{####1}}}%
  \@namedef{PYG@tok@kt}{\def\PYG@tc####1{\textcolor[HTML]{8B4E8F}{####1}}}%
  \@namedef{PYG@tok@ow}{\def\PYG@tc####1{\textcolor[HTML]{8B4E8F}{####1}}}%
  \@namedef{PYG@tok@s}{\def\PYG@tc####1{\textcolor[HTML]{4F8A55}{####1}}}%
  \@namedef{PYG@tok@s1}{\def\PYG@tc####1{\textcolor[HTML]{4F8A55}{####1}}}%
  \@namedef{PYG@tok@s2}{\def\PYG@tc####1{\textcolor[HTML]{4F8A55}{####1}}}%
  \@namedef{PYG@tok@sa}{\def\PYG@tc####1{\textcolor[HTML]{4F8A55}{####1}}}%
  \@namedef{PYG@tok@sd}{\def\PYG@tc####1{\textcolor[HTML]{4F8A55}{####1}}}%
  \@namedef{PYG@tok@si}{\def\PYG@tc####1{\textcolor[HTML]{4F8A55}{####1}}}%
  \@namedef{PYG@tok@m}{\def\PYG@tc####1{\textcolor[HTML]{A06D20}{####1}}}%
  \@namedef{PYG@tok@mi}{\def\PYG@tc####1{\textcolor[HTML]{A06D20}{####1}}}%
  \@namedef{PYG@tok@mf}{\def\PYG@tc####1{\textcolor[HTML]{A06D20}{####1}}}%
  \@namedef{PYG@tok@mh}{\def\PYG@tc####1{\textcolor[HTML]{A06D20}{####1}}}%
  \@namedef{PYG@tok@mo}{\def\PYG@tc####1{\textcolor[HTML]{A06D20}{####1}}}%
  \@namedef{PYG@tok@nb}{\def\PYG@tc####1{\textcolor[HTML]{9A702D}{####1}}}%
  \@namedef{PYG@tok@c}{\def\PYG@tc####1{\textcolor[HTML]{7A887E}{####1}}}%
  \@namedef{PYG@tok@c1}{\def\PYG@tc####1{\textcolor[HTML]{7A887E}{####1}}}%
  \@namedef{PYG@tok@cm}{\def\PYG@tc####1{\textcolor[HTML]{7A887E}{####1}}}%
  \@namedef{PYG@tok@cp}{\def\PYG@tc####1{\textcolor[HTML]{6F8175}{####1}}}%
}
% 用 fvextra 的全宽行盒绘制白/极浅灰逻辑行底纹。breaklines 产生的
% 视觉续行仍位于同一个 parbox 中，因此保持与源代码逻辑行相同的底色。
\newcommand{\AppendixFullWidthStripe}[2]{%
  \setlength{\FV@TmpLength}{\fboxsep}%
  \setlength{\fboxsep}{0pt}%
  \colorbox{#1}{%
    \setlength{\fboxsep}{\FV@TmpLength}%
    \rlap{\strut#2}%
    \hspace{\linewidth}%
    \ifx\FV@RightListFrame\relax\else
      \hspace{-\FV@FrameSep}\hspace{-\FV@FrameRule}%
    \fi
    \ifx\FV@LeftListFrame\relax\else
      \hspace{-\FV@FrameSep}\hspace{-\FV@FrameRule}%
    \fi
  }%
  \hss
}
\newcommand{\AppendixZebraLine}[1]{%
  \ifodd\value{FancyVerbLine}%
    \AppendixFullWidthStripe{AppendixStripeWhite}{#1}%
  \else
    \AppendixFullWidthStripe{AppendixStripeGray}{#1}%
  \fi
}
\AtBeginEnvironment{MintedVerbatim}{\AppendixPygmentsPalette}
\makeatother

% 透明、无左右框的可分页外壳。overlay 会在 unbroken/first/middle/last
% 每个物理分片上执行，从而显式重复绘制上下细横线。
\newtcolorbox{AppendixCodePageFrame}{%
  enhanced jigsaw,
  breakable,
  width=\linewidth,
  colback=white,
  colframe=AppendixHorizontalRule,
  boxrule=0pt,
  leftrule=0pt,
  rightrule=0pt,
  toprule=0pt,
  bottomrule=0pt,
  boxsep=0pt,
  left=0pt,
  right=0pt,
  % 为每个物理分页片段的横线保留极窄的无遮挡带，避免首/末行
  % 的全宽 zebra 色块覆盖横线主体。
  top=1.5pt,
  bottom=1.5pt,
  arc=0pt,
  outer arc=0pt,
  sharp corners,
  before skip=0.15\baselineskip,
  after skip=0.20\baselineskip,
  lines before break=3,
  overlay={%
    \draw[AppendixHorizontalRule,line width=0.5pt]
      ([yshift=-0.25pt]frame.north west)--([yshift=-0.25pt]frame.north east);%
    \draw[AppendixHorizontalRule,line width=0.5pt]
      ([yshift=0.25pt]frame.south west)--([yshift=0.25pt]frame.south east);%
  }
}

\setminted{
  encoding=utf8,
  linenos,
  numberblanklines=true,
  numbersep=6pt,
  breaklines,
  breakanywhere,
  autogobble,
  tabsize=4,
  obeytabs=true,
  fontfamily=myFont,
  fontsize=\fontsize{10.2pt}{13.5pt}\selectfont,
  baselinestretch=1.0,
  frame=leftline,
  framesep=5pt,
  framerule=0.55pt,
  rulecolor=\color{AppendixDivider},
  xleftmargin=3.15em,
  xrightmargin=0pt,
  listparameters={\setlength{\topsep}{0pt}\setlength{\partopsep}{0pt}\setlength{\parsep}{0pt}},
  resetmargins=true,
  style=lovelace
}

\newcommand{\AppendixStart}{%
  \section{附录}%
}

\newcommand{\AppendixPart}[2]{%
  \par\vspace{1.3\baselineskip}%
  \phantomsection
  \noindent{\songti\fontsize{12pt}{18pt}\selectfont\bfseries 附录 #1：#2}\par
  \addcontentsline{toc}{subsection}{附录 #1：#2}%
  \vspace{0.55\baselineskip}%
}

\newcommand{\AppendixAttachmentTableBegin}{%
  \begingroup
  \renewcommand{\arraystretch}{1.18}%
  \setlength{\tabcolsep}{4pt}%
  \small
  \begin{longtable}{>{\centering\arraybackslash}p{1.0cm}p{3.0cm}p{10.0cm}}
  \toprule
  序号 & 文件夹名称 & 主要内容 \\
  \midrule
  \endfirsthead
  \toprule
  序号 & 文件夹名称 & 主要内容 \\
  \midrule
  \endhead
  \midrule
  \multicolumn{3}{r}{续下页}\\
  \endfoot
  \bottomrule
  \endlastfoot
}

\newcommand{\AppendixAttachmentRow}[3]{%
  #1 & #2 & #3 \\
}

\newcommand{\AppendixAttachmentTableEnd}{%
  \end{longtable}
  \endgroup
}

% #1: 可选 minted 参数（支持 firstline/lastline/firstnumber）
% #2: 代码片段标题；#3: Pygments lexer；#4: 真实源码路径。
\newcommand{\AppendixCodeFile}[4][]{%
  \Needspace{6\baselineskip}%
  \par\vspace{0.35\baselineskip}%
  \noindent{\color{AppendixCaption}\songti\fontsize{10.5pt}{15pt}\selectfont #2}\par
  \vspace{0.15\baselineskip}%
  \begingroup
  \renewcommand{\theFancyVerbLine}{\AppendixCodeLatin\fontsize{8.6pt}{13.5pt}\selectfont\color{AppendixLineNumber}\arabic{FancyVerbLine}}%
  \let\FancyVerbHighlightLine\AppendixZebraLine
  \begin{AppendixCodePageFrame}%
    \inputminted[#1]{#3}{#4}%
  \end{AppendixCodePageFrame}%
  \endgroup
  \par\vspace{0.25\baselineskip}%
}

\newcommand{\AppendixEnd}{\par}
